\section{Graphs}
\label{sec:graphs}

Every section so far has computed with objects that carry a \emph{value}---an
integer, a polynomial, a matrix of numbers. A graph is different in kind: it is
pure \emph{structure}, a set of things together with a record of which of them
are joined, and nothing else. Strip the labels off a road map, a molecule, a
citation network, or a circuit and what remains---who is adjacent to
whom---is a graph. The mathematics of these objects, \emph{graph theory}, asks
questions that have no analogue for a single number: is everything reachable
from everything else, what is the cheapest route between two points, how few
colours suffice to keep neighbours distinct?\index{graph}

Mathilda represents a graph with the same machinery it uses for everything
else. A graph is an ordinary expression---no new kind of value is introduced,
just a head \B{Graph} wrapping a list of vertices and a list of edges---so the
generic tools you already know (\B{Part}, \B{Map}, \B{ReplaceAll}) reach inside
it, and its matrix views drop straight into the linear algebra of
\S\ref{sec:linalg}. We build graphs and read them back; take their matrix views
and hand them to \B{Eigenvalues}; call the ready-made generators; measure local
structure with vertex degrees; compute shortest paths---first by hop count, then
by weight, which is Dijkstra's algorithm; and finally colour a graph with the
fewest colours possible, a problem easy to state and, in the exact, NP-hard to
solve.

\subsection{Building a graph}
\label{ssec:graph-construct}

A graph is written \B{Graph}\mcode{[v, e]}, where \mcode{v} is the list of
vertices and \mcode{e} the list of edges. An edge is a pair of vertices with a
head that records whether it has a direction: \mcode{DirectedEdge[u, v]} for an
arrow \(u \to v\), \mcode{UndirectedEdge[u, v]} for a plain link. You rarely
type those heads, because the two rule arrows are accepted as
shorthand---\mcode{u -> v} (a \B{Rule}) becomes a directed edge and
\mcode{u <-> v} (the \mcode{TwoWayRule} operator) an undirected
one.\index{graph!edge}\index{graph!directed}\index{graph!undirected} Vertices
may be any expressions at all: integers here, but symbols, strings, or lists
work just as well.

\begin{codepairs}
\pair{graphs/construction/1}{A directed four-cycle. In standard output a graph
prints as a terse summary---its vertex and edge counts---not its full guts.}
\pair{graphs/construction/2}{\B{InputForm} shows the canonical form, which
round-trips through the parser: the arrows survive as written.}
\pair{graphs/construction/3}{\B{VertexList} reads back the vertices in canonical
order\ldots}
\pair{graphs/construction/4}{\ldots and \B{EdgeList} the edges, in the directed
form the shorthand produced.}
\pair{graphs/construction/5}{\B{EdgeCount} is one of the cardinalities
(\B{VertexCount} is the other).}
\end{codepairs}

You need not spell out the vertices. Given only the edges, \B{Graph} derives the
vertex set from them, in order of first appearance; and undirected edges print
back with the \mcode{<->} arrow:

\begin{codepairs}
\pair{graphs/construction/6}{With no vertex list, the vertices are read off the
edges---here \(a, b, c\).}
\pair{graphs/construction/7}{An undirected path, its edges echoed in
\mcode{UndirectedEdge} shorthand.}
\pair{graphs/construction/8}{Construction \emph{validates}: this self-loop is
rejected, so \B{GraphQ} reports the result is not a well-formed graph.}
\end{codepairs}

The last line is worth dwelling on. Mathilda's graphs are \emph{simple}: no
self-loops, no parallel edges, no vertex named in an edge but missing from the
vertex list. Malformed input is not silently repaired---it is left unevaluated,
and \B{GraphQ} is the predicate that tells you whether you hold a valid graph.

\usagebox{Graph}

\calloutnote{Under the hood}{There is no new \texttt{EXPR} tag for graphs. A
graph is literally the expression \texttt{Graph[List[v1, ...], List[e1, ...]]},
which is why \B{Part} and \B{Map} and pattern replacement
work on it unchanged---the uniformity principle of \S\ref{sec:session-state} at
work. On construction, \texttt{src/graph/construct.c} normalizes the edge
shorthand to \texttt{DirectedEdge}/\texttt{UndirectedEdge}, then validates
simplicity and endpoint membership before handing back the canonical form. The
terse \texttt{Graph[<n vertices, m edges>]} you see is only how the printer
summarizes it; \B{FullForm} exposes the tree.}

\subsection{Matrix views}
\label{ssec:graph-matrix}

The oldest bridge between graph theory and computation is the \emph{adjacency
matrix}: the \(n \times n\) array whose \((i,j)\) entry is \(1\) when there is an
edge from vertex \(i\) to vertex \(j\) and \(0\) otherwise.\index{adjacency
matrix} \B{AdjacencyMatrix} builds it, and---this is the point---it returns an
ordinary matrix, a list of lists, so every tool of linear algebra applies to it
directly. Take the directed four-cycle \(g\) from above:

\begin{codepairs}
\pair{graphs/matrices/2}{The adjacency matrix of the directed 4-cycle: a single
\(1\) in each row, walking the cycle.}
\pair{graphs/matrices/3}{\B{Det} of it is just \(\pm1\)---the cycle is a
permutation matrix.}
\pair{graphs/matrices/4}{The \emph{undirected} 4-cycle is symmetric, so
\B{Eigenvalues} are real: its spectrum \(\{-2, 2, 0, 0\}\).}
\end{codepairs}

That last computation is a first taste of \emph{spectral graph
theory}\index{spectral graph theory}: the eigenvalues of the adjacency matrix
encode structural facts about the graph---connectivity, bipartiteness, the count
of closed walks---and here the extreme values \(\pm2\) are the signature of a
2-regular bipartite cycle. Because the matrix is a genuine \B{List} of lists,
nothing about the graph subsystem had to be taught linear algebra; the two meet
on common ground.

Edges may also carry \emph{weights}---a distance, a cost, a capacity---supplied
as a third argument \mcode{EdgeWeight -> \{w1, ...\}}, matched to the edges by
position. \B{EdgeWeight} reads them back, and \B{WeightedAdjacencyMatrix} places
each edge's weight where the plain adjacency matrix would place a \(1\):

\begin{codepairs}
\pair{graphs/matrices/6}{The weights, in edge order, of a two-edge weighted
graph \(w\).}
\pair{graphs/matrices/7}{Its weighted adjacency matrix carries \(5\) and \(7\)
in place of ones.}
\pair{graphs/matrices/8}{An unweighted graph defaults every weight to \(1\), so
its weighted matrix coincides with its plain one.}
\end{codepairs}

\calloutnote{Under the hood}{\B{AdjacencyMatrix} (\texttt{src/graph/adjmat.c})
materializes a dense \(n\times n\) matrix in canonical vertex order, so
\B{Det}, \B{Tr}, and \B{Eigenvalues} consume it with no conversion.
Its inverse, \B{AdjacencyGraph}, rebuilds a graph from a 0/1 matrix, and
\B{IncidenceMatrix} gives the \(|V|\times|E|\) vertex--edge view instead. For a
large sparse graph the dense matrix is the wrong representation---the search
routines below never form it, walking an integer-indexed adjacency list
instead---but for the small graphs where you actually want a determinant or a
spectrum, dense is exactly right.}

\subsection{Ready-made graphs}
\label{ssec:graph-generators}

Named families of graphs recur so often that Mathilda generates them directly,
each on the vertices \(1,\dots,n\). \B{StarGraph}\mcode{[n]} joins a central hub
to \(n-1\) leaves---the simplest hub-and-spoke topology, and a graph you will
meet again as a tree.\index{star graph}

\begin{codepairs}
\pair{graphs/generators/1}{The star on five vertices: a hub and four leaves,
four edges in all.}
\pair{graphs/generators/2}{Its edges all emanate from vertex \(1\).}
\pair{graphs/generators/3}{So vertex \(1\) has degree \(4\) and every leaf
degree \(1\)---the defining shape of a star.}
\end{codepairs}

For experiments and average-case behaviour you want a graph drawn \emph{at
random}. \B{RandomGraph}\mcode{[\{n, m\}]} samples uniformly among all graphs with
\(n\) vertices and exactly \(m\) edges---the \(G(n,m)\) model of Erd\H{o}s and
R\'enyi.\index{random graph} It draws from the same seeded stream as the other
random builtins, so \B{SeedRandom} makes any run reproducible---which is exactly
why every example in this book can be regenerated:

\begin{codepairs}
\pair{graphs/generators/5}{With the stream seeded, a random graph on six
vertices with eight edges---the same one on every run.}
\pair{graphs/generators/6}{Its eight edges, drawn uniformly without
replacement.}
\pair{graphs/generators/7}{A second argument asks for a list of independent
samples; here three of them.}
\end{codepairs}

\calloutnote{Under the hood}{\B{RandomGraph} (\texttt{src/graph/generators.c})
enumerates the \(\binom{n}{2}\) candidate edges and hands them to
\B{RandomSample}, so it shares the one global RNG that \B{SeedRandom} controls
and \B{RandomReal} draws from. The count form \mcode{RandomGraph[\{n,m\}, k]}
costs a full candidate materialization per sample, so its time and memory scale
as \(O(k\,n^2)\)---fine for the small graphs of a teaching example, a real
consideration at scale. \B{CompleteGraph}, \B{CycleGraph}, and \B{PathGraph}
round out the family.}

\subsection{Local structure: vertex degrees}
\label{ssec:graph-degree}

The most local fact about a vertex is its \emph{degree}: how many edges touch
it. In a directed graph this splits in two. The \emph{out-degree} counts edges
leaving a vertex, the \emph{in-degree} edges arriving---and the distinction is
the whole content of ``who points at whom''.\index{degree (graph)!in-degree}\index{degree
(graph)!out-degree} Take a directed graph with one extra chord:

\begin{codepairs}
\pair{graphs/degrees/2}{The 4-cycle \(1\!\to\!2\!\to\!3\!\to\!4\!\to\!1\) with an
added chord \(1\!\to\!3\).}
\pair{graphs/degrees/3}{\B{VertexInDegree}: vertex \(3\) is now pointed at twice
(by \(2\) and by \(1\)), the rest once.}
\pair{graphs/degrees/4}{\B{VertexOutDegree}: vertex \(1\) now emits twice, the
rest once.}
\pair{graphs/degrees/5}{A second argument asks about a single vertex---the
in-degree of \(3\).}
\pair{graphs/degrees/6}{And the out-degree of \(1\).}
\end{codepairs}

The two degree sequences are not independent: every directed edge contributes
exactly one to some vertex's out-degree and one to another's in-degree, so their
sums are equal, each counting the edges once. That is the directed form of the
\emph{handshaking lemma}, and summing either sequence recovers
\B{EdgeCount}. For an undirected graph the plain \B{VertexDegree} counts incident
edges, as it did for the star above.

\subsection{Shortest paths}
\label{ssec:graph-shortest}

Now a global question: what is the shortest route from one vertex to another?
The answer depends on what ``short'' means. With no weights, length is
\emph{hop count}---the number of edges---and the shortest path is found by
breadth-first search. With weights, length is the \emph{total weight} along the
path, and finding the cheapest route is the province of Dijkstra's
algorithm.\index{shortest path} \B{FindShortestPath} returns the route as a
vertex list, \B{GraphDistance} its length; both make the same choice
automatically, keyed on whether the graph carries usable edge weights.

Consider a small road network among five towns, its edge weights the distances
between them. Let \mcode{roads} be that weighted graph and \mcode{plain} the
same towns and roads with the distances forgotten:

\begin{codepairs}
\pair{graphs/shortestpath/4}{The cheapest route from \(a\) to \(e\) threads
through all five towns\ldots}
\pair{graphs/shortestpath/5}{\ldots for a total distance of \(8\).}
\pair{graphs/shortestpath/7}{Forget the weights and the shortest path is the one
with \emph{fewest hops}: straight across, \(a\!\to\!c\!\to\!e\).}
\pair{graphs/shortestpath/8}{That route is only two edges long\ldots}
\end{codepairs}

Here is the lesson of weighting in one pair of answers. The two-hop route
\(a\to c\to e\) looks shortest, and by hop count it is; but its edges are heavy
(\(5\) and \(8\), total \(13\)), while the four-hop detour
\(a\to b\to c\to d\to e\) stays on cheap roads and totals only \(8\). The fewest
steps and the least distance need not agree, and Dijkstra finds the latter.
Because the weights were exact integers, so is the distance---Mathilda does not
launder an exact answer through floating point:

\begin{codepairs}
\pair{graphs/shortestpath/9}{Rational weights stay rational: the distance
\(1/2 + 1/3 = 5/6\) is returned exactly, not as a decimal.}
\end{codepairs}

\calloutnote{Theory}{\emph{Dijkstra's algorithm}\index{Dijkstra's algorithm}
grows a set of vertices whose shortest distance from the source is settled,
repeatedly admitting the unsettled vertex of least tentative distance and
relaxing its neighbours. Its correctness rests on one hypothesis:
\emph{non-negative weights}. If every edge is non-negative, then the moment a
vertex is admitted its tentative distance cannot be beaten by any later route,
because any such route would have to pass through an as-yet-farther vertex and
so be no shorter. Breadth-first search is exactly Dijkstra with every weight
equal to \(1\): the frontier expands one hop at a time, settling vertices in
order of hop count. So the unweighted and weighted cases are one algorithm, and
Mathilda dispatches between them on the weights alone.}

\calloutnote{Under the hood}{Both routines (\texttt{src/graph/shortestpath.c})
build an integer-indexed adjacency on demand---never the dense matrix of
\S\ref{ssec:graph-matrix}---and run a plain \(O(V^2)\) Dijkstra, no priority
queue, which is the right trade for the small graphs the subsystem targets. The
floating-point accumulator inside Dijkstra only \emph{orders} the frontier; the
returned \B{GraphDistance} is recomputed exactly along the chosen path, so an
integer- or rational-weighted graph yields an exact answer. A \emph{negative}
weight would break Dijkstra's hypothesis, so rather than return a wrong answer
the routines fall back to unweighted breadth-first search when a weight is
negative, symbolic, or absent---there is no Bellman--Ford path here.}

\usagebox{FindShortestPath}

\calloutnote{Pitfall}{Because a negative or symbolic weight silently demotes the
computation to hop-count BFS, a graph you \emph{intended} to be weighted but
whose weights are not all non-negative numbers will quietly answer the wrong
question---fewest hops, not least cost. When weights matter, check that they are
all present and non-negative; \B{GraphDistance} returning an integer where you
expected a weighted total is the tell.}

\subsection{Colouring a graph}
\label{ssec:graph-colouring}

The last problem is deceptively simple to state. A \emph{proper colouring}
assigns a colour to each vertex so that no edge joins two vertices of the same
colour; the least number of colours that admits a proper colouring is the
graph's \emph{chromatic number} \(\chi(G)\).\index{chromatic number}\index{graph
colouring} \B{FindVertexColoring} returns a colouring as a list of integer
labels in vertex order, and---crucially---it returns a \emph{minimal} one: the
largest label it uses \emph{is} \(\chi(G)\). Three graphs tell the story:

\begin{codepairs}
\pair{graphs/coloring/1}{An even cycle is \emph{bipartite}: two colours suffice,
alternating around the ring.}
\pair{graphs/coloring/2}{An \emph{odd} cycle cannot be 2-coloured---close the
ring and the two ends clash---so it needs three.}
\pair{graphs/coloring/3}{In the complete graph every pair is adjacent, so no two
vertices may share a colour: \(\chi(K_4) = 4\).}
\pair{graphs/coloring/4}{Reading off the chromatic number is just the largest
label used.}
\end{codepairs}

These three cases map out the extremes. A graph is 2-colourable exactly when it
is bipartite---when it has no odd cycle---so even cycles, paths, trees, and the
star all take two colours, while a single odd cycle already forces three. At the
other end, the complete graph \(K_n\) forces all \(n\): every vertex neighbours
every other, so all must differ. The star makes the bipartite point cleanly:

\begin{codepairs}
\pair{graphs/coloring/5}{The star is bipartite (hub against leaves), so two
colours suffice however many leaves there are.}
\end{codepairs}

\calloutnote{Theory}{That \(\chi(G)=2\) iff \(G\) is bipartite iff \(G\) has no
odd cycle is a theorem with a one-line proof: 2-colouring \emph{is} a
2-partition with no edge inside a part, and an odd cycle admits no such
partition. But for \(k\ge 3\), deciding whether \(\chi(G)\le k\) is
\emph{NP-complete}\index{NP-hardness}---one of Karp's original 21 problems---so
no method is known that computes \(\chi(G)\) in polynomial time in general. This
is why \B{FindVertexColoring} carries a warning the shortest-path functions do
not: computing it exactly can be expensive.}

\calloutnote{Under the hood}{A greedy colouring is cheap but not minimal;
Mathilda's \B{FindVertexColoring} (\texttt{src/graph/vertexcoloring.c}) instead
proves minimality by exact search---a DSATUR branch-and-bound bracketed between
a cheap DSATUR upper bound and a clique-based lower bound---and this is Wolfram's
own \texttt{"BacktrackingDS"} method. Because the problem is NP-hard, the search
refuses rather than lie: above \(128\) vertices, or after an eight-million-node
budget, it returns \emph{unevaluated}---never a valid-but-larger colouring. Cost
is driven by edge \emph{density}, not vertex count, so a sparse graph on many
vertices colours instantly while a dense small one can exhaust the budget. To
bound the time, wrap the call in \B{TimeConstrained}.}

\usagebox{FindVertexColoring}

\subsection{Drawing a graph}
\label{ssec:graph-plot}

A picture often says what a vertex list cannot. \B{GraphPlot}\mcode{[g]} lays the
vertices out and draws the edges, producing a \texttt{Graphics} object that
renders through the same pipeline as \B{Plot} and \B{Show} (a window when
graphics support is compiled in, a text placeholder otherwise). Being a picture,
it has no printed transcript to reproduce here---it is output for the eye, not
the page---so we only note its existence; the graphics chapter
(Chapter~\ref{ch:graphics}) is where drawing is taken up in earnest.

\subsection*{Where this connects}

Graphs are where Mathilda's ``everything is an expression'' principle earns its
keep: a graph is a plain \B{Graph} tree, so the generic structural tools operate
on it and its matrix views are ordinary matrices. That last fact ties this
section to the linear algebra of \S\ref{sec:linalg}---the adjacency matrix is
the object spectral graph theory studies, and \B{Eigenvalues}, \B{Det}, and
\B{Dot} apply to it with nothing in between. The reproducible sampling of
\B{RandomGraph} rests on the same seeded random-number stream as the simulations
of \S\ref{ssec:stat-distributions}. And the two search algorithms met here are
templates for a wider family---the same integer-indexed traversal underlies
\B{ConnectedComponents}, \B{FindSpanningTree}, and \B{VertexConnectivity}---while
the exact, refuse-rather-than-approximate discipline of \B{FindVertexColoring}
is the same stance Mathilda takes throughout: an honest ``I cannot'' in place of
a plausible wrong answer. For the full options of any function named here, its
reference page is one \texttt{?Name} away at the prompt.
